% =====================================================================
%  PAPER 2 (working skeleton): It from Bit via G\"odel, II ---
%  Spacetime from the FANOUT Boundary
%  CONVERTED 2026-07-07 from paper2_metric_identification_2026-07-07.md
%  (the seed remains the idea-log; this file is the manuscript).
%
%  STATE: sections 2, 4, 5, 6 are drafted at paper register; sections
%  3, 7, 8 are structured stubs with their content minted as numbered
%  open problems. Six problems minted; the paper opens, like its
%  parent, with its debts on display.
%
%  PREAMBLE: mirrors Paper 1's machinery (amsthm + thmtools
%  openproblem + the Open Problems list). SYNC NOTE for JS: copy the
%  exact \declaretheorem styling from Paper 1's main.tex for visual
%  identity, and point \addbibresource at the shared library --- every
%  Paper-1 key (doraumuch2026relative, brillouin1956science,
%  bose2017spin, marlettovedral2017gravitationally,
%  procopio2015experimental, toffoli_2003, Jacobson1995, Penrose1996,
%  small_2006, Landauer_1961, zurek_2009 ...) is reused here; four new
%  keys are in paper2_starter_refs_2026-07-07.ris.
%
%  CROSS-PAPER REFERENCES: \Pone{...} marks references into Paper 1
%  (its sections and open problems); resolve to a self-citation key at
%  Paper 1's publication. British English; \autocite; no bullets in
%  prose (one table, deliberate).
% =====================================================================
\documentclass[11pt,a4paper]{article}
\usepackage[T1]{fontenc}
\usepackage{amsmath, amssymb, amsthm}
\usepackage{thmtools}
\usepackage[british]{babel}
\usepackage{csquotes}
\usepackage[backend=biber, style=numeric-comp, sorting=none]{biblatex}
\addbibresource{self-ref-2026.bib} % shared library; TODO(JS): confirm path
\newtheorem{proposition}{Proposition}
\newtheorem{conjecture}{Conjecture}
\declaretheorem[style=definition, name=Open Problem]{openproblem}
\newcommand{\docversion}{Draft 0.1 --- 2026-07-07}
\newcommand{\Pone}[1]{[Paper~I, #1]} % TODO: resolve at publication
\newcommand{\fanout}{\textsc{fanout}}
\title{It from Bit via G\"odel, II:\\ Spacetime from the \fanout{} Boundary}
\author{John Small}
\date{\docversion}

\begin{document}
\maketitle

\begin{abstract}
% TODO(JS): working abstract; replace at will.
The first paper of this programme derived the structure of quantum mechanics and the Standard
Model from self-reference, and defined a spacetime interval by counting: computable steps counted
real, non-computable bits counted imaginary. This paper identifies that computational metric with
the spacetime of general relativity. The identification has a kinematic half --- the
Margolus--Levitin theorem and the de~Broglie relation make energy--momentum and step-count density
the same data in two currencies, with $\hbar$ the exchange rate --- and a dynamic half, for which
the trail is now paved: Dorau and Much have proved, with modular-theoretic rigour, that the
semiclassical Einstein equations follow from quantum relative entropy once the entropy--area law
is assumed, and the entropy--area law is exactly what this programme's counting supplies. We state
the count-density dictionary as a conjecture, mint the problems that would close it, give the
four-constant junction by which one \fanout{} locking is billed simultaneously in $c$, $\hbar$,
$G$ and $k_{B}$, and supply the two-level ontology under which thermodynamic derivations of the
Einstein equations stop consuming what they produce. Gravity, we argue, lives in the
post-\fanout{} regime: the quantum layer beneath it is pre-geometric, and the experimental
programme this reading organises --- gravitationally induced entanglement predicted without a
quantised metric, standard decoherence rates, no recorded indefinite causal order --- is already
under construction in the laboratories.
\end{abstract}

\section{Introduction and programme}
\label{sec:p2-intro}

% TODO(JS): full introduction at the final-block stage, as in Paper 1.
% The paragraphs below are the load-bearing framing, drafted.

The first paper of this programme \Pone{} ended its spacetime section with a deferral: the deepest
questions --- why the counting metric should be \emph{the} metric, and what makes its geometry
obey Einstein's equations --- were owed to a second paper. This is that paper. Its central claim
can be put in one sentence: the metric general relativity builds from energy and momentum and the
metric the first paper builds from computable and non-computable counts are the same metric,
because energy--momentum and count-density are the same data in two currencies. The claim has a
provenance: the counting convention, spacelike distance in imaginary units of information, was
proposed in the author's CASYS paper two decades ago with its proof flagged, already then, as a
key open problem \autocite{small_2006}; the machinery assembled since --- the \fanout{} boundary,
the relational reading of causal structure, and now the modular-theoretic derivation of the
semiclassical Einstein equations by Dorau and Much \autocite{doraumuch2026relative} --- is what
lets the flag be answered rather than waved.

One remark orients the whole paper. The oldest tool of the programme --- Lawvere's fixed-point
shape, the inside description constraining what it describes --- turns out to be what the newest
and most rigorous result in the emergent-gravity literature is missing: read without a two-level
ontology, the thermodynamic derivations of Einstein's equations consume the arena they produce;
read with one, they are exactly the inside view deriving the law its own arena must satisfy. The
first paper's first move fixes the field's latest theorem. Sections~\ref{sec:p2-rosetta}
and~\ref{sec:p2-kinematic} state the identification and its kinematic gate, with
Section~\ref{sec:p2-energy} spending the dictionary's first dividend --- the non-localisability
of gravitational energy dissolved rather than defined around;
Section~\ref{sec:p2-dynamic} walks the paved trail; Section~\ref{sec:p2-junction} gives the
four-constant junction; Section~\ref{sec:p2-ontology} supplies the ontology;
Section~\ref{sec:p2-postfanout} states where gravity lives and what the laboratories can do about
it; Section~\ref{sec:p2-cognates} maps the neighbours. The problems are numbered where they arise
and indexed at the end, in the parent paper's habit: the debts are the invitation.

\section{The Rosetta line}
\label{sec:p2-rosetta}

The Margolus--Levitin theorem bounds the rate at which a system of mean energy $E$ above its
ground state can pass through orthogonal states: at most $2E/h$ distinguishable steps per unit
time \autocite{margoluslevitin1998}. Its spatial sibling is de~Broglie read computationally: a
system of momentum $p$ passes through orthogonal position states at rate $p/\hbar$ per unit
length. Together they say that, up to $\hbar$, the pair $(E, p)$ \emph{is} the pair (computable
steps per unit time, distinguishable marks per unit length). General relativity sources the metric
on $T_{\mu\nu}$; the first paper builds the metric from counting computational events; the
Margolus--Levitin--de~Broglie pair says the density of computational events \emph{is}
$T_{\mu\nu}$. The two definitions of the metric were never rivals to be reconciled: they are one
definition read through Planck's constant. There is a structural echo worth recording: the
temporal count rides $E$ and the spatial count rides $p$, so the real/imaginary split of the
computational distance maps onto the time/space block structure of the stress tensor exactly where
it should.

\begin{conjecture}[The count-density dictionary]\label{conj:dictionary}
Let the computational metric of \Pone{sec:space-of-computations} be a genuine pseudo-metric with
Minkowski signature (the content of \Pone{op:comp-metric}). Then the map sending a region's
computable-step density and distinguishable-mark density to $(E/\hbar, p/\hbar)$ identifies the
counting interval with the Lorentzian interval, and the stress--energy tensor with the density of
computational events, exactly and with $\hbar$ as the sole conversion.
\end{conjecture}

The conjecture is conditional by construction --- it inherits the parent paper's open problem as
its hypothesis --- and honest about its two gaps: the $\hbar$ bookkeeping must be stated with the
factor-of-two conventions of the Margolus--Levitin bound handled explicitly, and the counted
``step'' must be an \emph{orthogonalising} transition for the bound to bind at all, which the
\fanout{} images guarantee by construction and the definition must be made to say.
% TODO: the hbar bookkeeping paragraph (seed item 1).

\begin{openproblem}[The stress components]\label{op:p2-pressure}
Extend the dictionary beyond the $(E, p)$ pair: pressure and shear are momentum flux, and the
computational reading of flux --- distinguishable marks carried across a cut per unit time ---
must be defined and shown to reproduce the perfect-fluid form of $T_{\mu\nu}$ in the appropriate
regime. The wire-crossing bookkeeping of the relational network is the candidate machinery, and the
flagship instance is Isaacson's effective stress of the gravitational wave itself
(Section~\ref{sec:p2-energy}): re-bookkeeping activity as event-density after averaging.
\end{openproblem}

\section{Energy, undefined exactly where it should be}
\label{sec:p2-energy}

The dictionary of Section~\ref{sec:p2-rosetta} has a corollary that dissolves one of general
relativity's oldest embarrassments. Energy, on the counting reading, is a rate: computable steps
per unit time, information moving through time as momentum is information moving through space.
A rate needs a direction to be counted along, and in general relativity the time direction is not
an invariant of the theory but a local choice --- the light cone tips, timelike converts to
spacelike, and time translation belongs to the diffeomorphism gauge group rather than to any rigid
symmetry. The famous non-localisability of gravitational energy is therefore not a technical
misfortune awaiting a cleverer definition: it is what the dictionary \emph{predicts}. Energy is
exactly as well-defined as the temporal chart is coherent, which is why the textbook criterion ---
conserved energy if and only if a timelike Killing field exists --- translates verbatim: a Killing
time is a globally locked step-direction, and stationarity is temporal lock-uniformity, the chart
of Open Problem~\ref{op:p2-chart} extended in time.

Aoki's recent analysis gives the group-theoretic face of the same dissolution
\autocite{aoki2025energy}: attempts to restore conservation by assigning energy to the
gravitational field fail by Noether's second theorem \autocite{noether1918invariante} --- the
candidate currents are conserved identically, off-shell, so the pseudotensors and quasilocal
energies of the textbooks are constraints rather than dynamical charges, fragile under total
derivatives and vanishing in closed universes. The ledger reading says why this had to be so:
geometry does not carry counts because geometry \emph{is} the bookkeeping of counts, and asking
the gravitational field to have energy is asking the ledger to book itself --- the diagonal move
this programme polices everywhere, met here in its energetic clothing. What survives, in the same
analysis, is the framework's own currency: the conserved charge that replaces energy is the
streamline count, and for perfect fluids it is the entropy current
\autocite{aokihidakakawanashimada2025}. Curved spacetime conserves not the chart-relative
count-rate but the chart-free count of entries --- bits, worldlines, locked identities --- and the
prohibition of vacuum-to-matter transitions their charge enforces reads, in this vocabulary, as
``no entries appear in the ledger without formation events''.

The objection every reader will now raise deserves its answer in place: gravitational radiation
carries energy --- binary pulsars spin down at the quadrupole rate, interferometer mirrors move,
and the loudest events close their books short by solar masses. If there is no gravitational
energy, how? The verdict denies none of it; it denies a \emph{local, tensorial, chart-independent
energy density}, and the mature literature's actual inventory is precisely what the dictionary
predicts a tipping-bulk theory should possess. Energy exists at the matter endpoints, which are
locked systems; at the asymptotic boundaries, where the Bondi mass-loss formula is theorem-grade
geometry with no pseudotensor in it \autocite{bondivanderburgmetzner1962} --- null infinity being
the ultimate locked chart, so that Bondi energy is a boundary-ledger difference; and in the wave
zone only after Isaacson's averaging over several wavelengths
\autocite{isaacson1968gravitational}, which the dictionary reads exactly: \emph{the averaging
scale is the chart re-emergence scale}, the minimum region over which enough locking pattern
accumulates for an effective step-direction, hence a count-rate, to exist again. Pointwise in the
bulk there is no energy because the count-direction is the very thing oscillating --- the wave is
the ledger updating, the phonon of the network doing matter-side work wherever it meets a locked
system --- and the pseudotensors' century of frame-dependence is the locked/unlocked split made
visible. Conservation survives as what it operationally always was: a balance between locked
bookends, matter loss against boundary flux, silent about a location in the unlocked between. One
demand falls out and is filed where it belongs: the wave's own gravitating --- Isaacson's
back-reaction --- requires that re-bookkeeping activity itself count as event-density after
averaging, the flagship instance of Open Problem~\ref{op:p2-pressure}.

\begin{openproblem}[Worldlines as locked identities]\label{op:p2-worldlines}
Aoki and collaborators note that the non-termination and non-emergence of the streamlines
underlying their conserved charge is plausible but unproven. Derive it from record persistence:
a worldline is a locked identity chain (the individuation of \Pone{}'s confinement analysis), and
its non-termination should follow from the same practical irreversibility that grades the
\fanout{} locking --- turning their assumption into this framework's theorem, and the contact
point into a contribution.
\end{openproblem}

\section{The kinematic rung}
\label{sec:p2-kinematic}

% STUB. Content: the gate is Paper I's op:comp-metric, restated here with its
% Paper-2 face --- the record-side chart. Seed items: the wedge/modular target.

The identification's kinematic half waits on one gate, and the gate now has a sharper target than
when the parent paper priced it. \Pone{op:comp-metric} asks that the paired distance be a genuine
pseudo-metric with derived signature; the dynamic rung below adds the demand that the resulting
causal structure furnish \emph{local wedges with modular structure} --- the bifurcate-horizon
scaffolding on which the rigorous derivations run. We record the strengthened form.

\begin{openproblem}[The record-side chart]\label{op:p2-chart}
Construct, from the relational network's locked sector, the local chart a community of
high-redundancy observers shares: show that it is locally flat (locking locally uniform), that its
causal wedges carry the modular structure the Bisognano--Wichmann setting requires, and that the
construction reduces, at one qubit, to the record-side data of \Pone{prop:f1-dictionary}. This is
the parent's open problem with its quality bar raised to modular-theoretic standard.
\end{openproblem}

\section{The dynamic rung: the paved trail}
\label{sec:p2-dynamic}

Jacobson derived the Einstein equations as a thermodynamic equation of state
\autocite{Jacobson1995}; Dorau and Much have now carried the derivation to modular-theoretic
rigour: the Araki--Uhlmann relative entropy between the vacuum and coherent excitations on a
bifurcate Killing horizon equals the energy flux across it, and, \emph{under the assumption of the
Bekenstein--Hawking entropy--area law}, the semiclassical Einstein equations follow
\autocite{doraumuch2026relative}. The interface with this programme could not be cleaner, and we
state it as the paper's division of labour: \emph{their hypothesis is our conjecture's
conclusion}. They prove that local horizon structure plus the entropy--area law yields the
Einstein equations; this programme proposes to derive exactly those two inputs --- the
entropy--area law from counting (Section~\ref{sec:p2-junction}) and the horizon structure from the
network (Open Problem~\ref{op:p2-chart}). The composition of their theorem with these two
deliveries is the full route from self-reference to the field equations. Their coherent-state
scope, moreover, is not a limitation to be transcended but a domain announcement: coherent states
are the record-friendly states, and Section~\ref{sec:p2-ontology} argues that the semiclassical
equations they obtain are the \emph{exact} law of the regime their states inhabit.

\section{The four-constant junction}
\label{sec:p2-junction}

In the native ledger --- time in steps, space in $i \times$ bits --- the four fundamental
constants are the four exchange rates into human units, and a single graded-\fanout{} locking is
billed in all four currencies at once. The copy is a step ($c$: one bit of spatial relation per
causal step; natively $c = 1$). It is an orthogonalising transition ($\hbar$: the
Margolus--Levitin pricing of Section~\ref{sec:p2-rosetta}). Its permanence is paid ($k_{B}$:
Landauer's bill of $k_{B} T \ln 2$ per locked bit \autocite{Landauer_1961}, with the practical
irreversibility supplied by redundancy \autocite{zurek_2009}). And the locked bit increments the
boundary ledger ($G$: one bit-pair of native area). Because spacelike distance is imaginary, a
spacelike area is the product of two imaginary lengths and is \emph{negative} in native units ---
the entropy--area law's sign is $i^{2}$, and Bekenstein--Hawking in native currency reads
$A = -S$, coefficient owed \autocite{brillouin1956science}. Newton's constant is then the memory,
kept in human units, of how large one bit of space is; its native value is unity, and gravity's
weakness is the report that the bit is small.

The junction can be displayed in one line. One locking at the recorder's Unruh temperature
satisfies
\begin{equation}
\frac{\delta Q}{T} \;=\; k_{B} \ln 2 \, \cdot\, \mathrm{d}N_{\mathrm{bits}}
\;=\; \frac{k_{B} c^{3}}{4 G \hbar}\, \mathrm{d}A ,
\end{equation}
each equality one exchange rate; fed to the Clausius relation, $\hbar$ and $k_{B}$ cancel and the
Einstein equations emerge in $G$ and $c$ alone. The cancellation is a standing self-check for
every derivation in this paper: \emph{no final geometric equation of the framework's may contain
$\hbar$ or $k_{B}$}. What the junction motivates it does not fix: the $\mathcal{O}(1)$
coefficient --- bits per bit-pair, equivalently the bit's size in Planck lengths --- is the
entire physical content of ``deriving $G$'', and it is triple-routed.

\begin{openproblem}[The coefficient]\label{op:p2-coefficient}
Derive the bits-per-bit-pair coefficient of the native area law $A = -S$ by any of its three
routes --- the one-bit-per-closed-curve pricing of \Pone{}, the Bekenstein normalisation, or the
Landauer--Unruh energy audit of a single locking --- and show the three agree. Agreement would
over-determine the coefficient and complete the derivation of $G$; the minimal loop's
circumference-versus-enclosed-area bookkeeping on a network is the known hard part, and the
Jacobson-lineage circularity warning applies: the coefficient must be earned from the framework's
own counting, not borrowed back from the law it derives.
\end{openproblem}

\begin{openproblem}[The Newtonian limit]\label{op:p2-newton}
Show that computational load slowing local clocks reproduces the $1/r$ potential in the
appropriate regime, checked against the Regge construction of Lloyd's computational-universe
programme \autocite{lloyd2005quantum} --- the first place the identification can fail
quantitatively, and therefore the first computation to run.
\end{openproblem}

\begin{openproblem}[Equivalence from blindness]\label{op:p2-equivalence}
The Margolus--Levitin bound is species-blind: only $E$ enters, never the substrate. Derive the
universality of free fall as the clock-rate theorem's indifference to what is computing ---
locating whatever auxiliary premise the derivation secretly needs --- so that the equivalence
principle joins the parent paper's collection of derived postulates. A second face of the same
statement is geometric: local flatness as locally \emph{uniform locking}
(Section~\ref{sec:p2-ontology}).
\end{openproblem}

\section{The ontology: what the fixed background is}
\label{sec:p2-ontology}

% Drafted from seed 3c --- the paper's philosophical spine.

Every derivation in the Jacobson lineage, the rigorous ones most visibly, runs on a fixed arena:
modular theory needs a vacuum, the vacuum needs a wedge algebra, the wedge needs a horizon, the
flow needs a Killing field --- all objects of a given, per-solution non-dynamical background. The
output is the law that makes the background dynamical. Read at one level, the derivation consumes
what it produces, and the standard bootstrap defence --- only local flatness is assumed, and the
derived theory's own equivalence principle guarantees it --- answers the technical circularity
while leaving the ontological one untouched: the local patch is still treated as given.

The two-level reading dissolves the problem by saying what the fixed background \emph{is}: the
post-\fanout{} description --- the arena as recorded, locked, and shared at high redundancy. It
looks pre-existing and non-dynamical for precisely the reason classical records look objective,
and it is not ontic but epistemic-made-rigid. The dynamical, relational layer is the nascent
locking of Section~\ref{sec:p2-junction}. Relocated, the rigorous derivations become the inside
view deriving the law its own arena must satisfy --- the parent paper's founding shape --- and two
corollaries arrive at once. Semiclassicality is exact: $\langle T_{\mu\nu}\rangle$ rather than
$\hat{T}_{\mu\nu}$ is the record-side shadow of the matter operator, so the semiclassical Einstein
equations are not an approximation awaiting transcendence but the exact law of the locked regime
--- and the equations' textbook failure mode, macroscopic superpositions sourcing gravity toward
the empty middle, marks the regime boundary rather than a pathology: the Page--Geilker experiment
\autocite{pagegeilker1981} sits directly on that boundary, and its result --- gravity follows the
resolution, not the average --- is what relation-formation-at-locking predicts for free. And local
flatness acquires its second face: every locked patch looks the same because locking is what
looking classical means, which is the geometric half of Open Problem~\ref{op:p2-equivalence}.

% TODO: one paragraph engaging the equation-of-state moral's history (Jacobson's
% own caveats; Padmanabhan's atoms of spacetime) --- seed 3c's defence-and-limit
% passage, expanded with citations at drafting.

\section{Gravity lives in the post-\fanout{} regime}
\label{sec:p2-postfanout}

The parent paper states the thesis \Pone{sec:no-quantised-gravity}: causal structure is the
bookkeeping of locked relations, so a superposition of causal structures is a category error, and
the quantum layer beneath gravity is pre-geometric --- amplitudes over relation-formation, not
over formed geometries --- with gravitons permitted as the phonons of the locking network and
forbidden as fundamental metric quanta. This section owes that thesis its engagement with the two
literatures it reorganises. The indefinite-causal-order programme --- process matrices, the
quantum switch --- becomes the nascent layer's native formalism: the realised switches
\autocite{procopio2015experimental} demonstrate sub-threshold indefiniteness that definitises at
locking, the time-delocalised-subsystems analyses showing the photonic realisations need no exotic
spacetime are allies rather than deflations, and the boundary --- \emph{no indefinite causal order
survives amplification into records} --- is this paper's theorem target in that formalism.
% TODO: Oreshkov--Costa--Brukner engagement paragraph; mint keys at drafting.
The gravitationally-induced-entanglement programme
\autocite{bose2017spin, marlettovedral2017gravitationally} is met head-on: the framework predicts
the entanglement, because the pre-geometric layer is quantum, while denying that a positive result
witnesses a superposed metric --- the mediator is nascent relational structure, non-classical in
exactly the sense the LOCC-mediator theorems require, without being a graviton.
% TODO: the LOCC-theorem engagement in detail.

The reading's experimental content fits in one table, and the exposure is the point.

\begin{center}
\begin{tabular}{ll}
\hline
Probe & This framework's posture \\
\hline
BMV-class entanglement & Positive predicted; witnesses the pre-geometric layer \\
Decoherence rates & Standard only; no gravity-specific (Penrose-rate) collapse \\
Indefinite causal order & Real below threshold; never amplified into records \\
Page--Geilker-class tests & Gravity follows the locking, not the average \\
\hline
\end{tabular}
\end{center}

\begin{openproblem}[The nascent regime, formalised]\label{op:p2-nascent}
State the low-redundancy regime in the relative-entropy toolkit the rigorous derivations use:
the grading as a data-processing statement (each copy a channel, monotonicity the
irreversibility), the Landauer bill as the relative-entropy inequality its modern proofs already
make it, and the Wigner's-friend phase recovery of \Pone{sec:wf-reconciliation} as the un-forming
of a relation before its bill is paid --- superposed geometry as low-$R$ relational structure,
made precise.
\end{openproblem}

\section{Cognates and differentiators}
\label{sec:p2-cognates}

% STUB, drafted skeleton. Seed §4 + the Dorau--Much reference trail.

Three neighbours bound the territory. Lloyd's computational-universe programme
\autocite{lloyd2005quantum} builds geometry from a generic computation's causal structure; this
programme's computation is \emph{specified} --- a derived particle content, a \fanout{} boundary,
a Standard Model attached --- so where Lloyd shows that some computation yields a geometry, the
claim here is that this computation yields this universe, predictions included. The causal-set
programme's sequential growth dynamics \autocite{rideoutsorkin2000} postulates that spacetime
accretes events under stochastic rules; this programme proposes the identity of the accretion
event --- a \fanout{} crossing its redundancy threshold --- and the currency of its
stochasticity, the thermodynamic ledger. And the modular-theoretic derivations
\autocite{doraumuch2026relative}, with their reference trail through Longo, Casini and Hollands,
supply the rigorous macroscopic face whose microscopic ledger and two-level ontology this paper
offers.
% TODO: expand the trail's citations at drafting.

\section{Scope, honesty, and what this paper does not claim}
\label{sec:p2-scope}

% Drafted from seed §5 + the standing flags.

Three cautions govern everything above. The counted step must be an orthogonalising transition or
the Margolus--Levitin bound does not bind; the \fanout{} images satisfy this by construction and
every definition in this paper is to be audited against it. The dissolution of
Section~\ref{sec:p2-ontology} is a structural reading until Open Problem~\ref{op:p2-chart}
delivers the record-side chart --- the gap is named and its filling proposed, not yet performed.
And the one-bit-per-closed-curve pricing meets the computational-power-of-closed-timelike-curves
literature at an interface this paper maps and does not cross.
% TODO: the Deutsch / Aaronson--Watrous interface paragraph.

\renewcommand{\listtheoremname}{Open Problems of Paper II}
\listoftheorems[ignoreall, show={openproblem}]

\printbibliography
\end{document}
